% !TEX root = ../main.tex

\documentclass[../main.tex]{subfiles}

\graphicspath{{\subfix{../images/}}}

\begin{document}

	\textbf{Актуальность исследования.} Современный человек сталкивается с большим количеством личных целей, задач и амбиций. Использование мобильных приложений и цифровых сервисов стало неотъемлемой частью процесса самосовершенствования и организации повседневной деятельности. В условиях растущей информационной нагрузки пользователи испытывают трудности в управлении целями, отслеживании прогресса и поддержании мотивации на протяжении длительного времени.

	Геймификация~--- применение игровых механик в неигровых контекстах~--- представляет собой активно развивающееся направление в области проектирования пользовательского опыта. Исследования последних лет демонстрируют значительный потенциал геймификации для повышения мотивации и вовлечённости пользователей. Системы достижений, визуализация прогресса и поэтапная декомпозиция целей способствуют формированию устойчивых поведенческих паттернов и поддержанию внутренней мотивации при достижении поставленных задач.

	Анализ существующих решений на рынке мобильных приложений для отслеживания достижений и привычек выявил ряд существенных недостатков: нестабильная работа интерфейса, отсутствие возможности обмена пользовательскими наборами достижений, ограниченная кроссплатформенность и зависимость от сетевого подключения. При этом потребность в специализированном инструменте, сочетающем геймифицированный подход к целеполаганию с возможностью совместного использования контента, сохраняет свою актуальность.

	Учитывая вышеперечисленное, было принято решение разработать кроссплатформенное приложение~--- \textit{Achi}, которое позволит пользователям структурировать цели в виде системы достижений, организованных в пакеты, отслеживать прогресс выполнения отдельных шагов и обмениваться наборами достижений с другими пользователями посредством уникальных кодов. Приложение должно обеспечивать синхронизацию данных с сервером, поддерживать различные типы прогресса (бинарный и инкрементальный) и предоставлять единый пользовательский опыт на платформах Android и~Web.

	\textbf{Новизна исследования} заключается в разработке кроссплатформенного клиентского приложения на базе Kotlin Multiplatform и Compose Multiplatform, реализующего иерархическую модель достижений (пакеты~$\rightarrow$~достижения~$\rightarrow$~шаги) с механизмом обмена пакетами по коду, оптимистичным обновлением прогресса и автоматической генерацией API-клиента из OpenAPI-спецификации. В отличие от рассмотренных аналогов, предлагаемое решение объединяет стабильный декларативный интерфейс, серверную синхронизацию прогресса и единую кодовую базу для нескольких целевых платформ.

	\textbf{Целью} данной работы является разработка кроссплатформенного мобильного приложения для трекинга достижений, применяющего принципы геймификации для повышения мотивации пользователей при достижении персональных целей.

	\textbf{Объектом исследования} являются процессы и средства разработки кроссплатформенных мобильных приложений в области персональной продуктивности с элементами геймификации.

	\textbf{Предметом исследования} являются методы и инструменты проектирования и реализации кроссплатформенного приложения для отслеживания достижений на базе Kotlin Multiplatform, включая выбор архитектурного паттерна, проектирование структур данных, интеграцию с серверным API и организацию пользовательского интерфейса.

	Для достижения поставленной цели необходимо решить следующие \textbf{задачи}:
	\begin{enumerate}[leftmargin=*]
		\item провести анализ предметной области и существующих решений в области трекинга достижений и привычек, выявить их функциональные и технические недостатки;
		\item сформулировать функциональные и нефункциональные требования к разрабатываемому приложению;
		\item провести анализ технологий кроссплатформенной разработки и обосновать выбор стека (Kotlin Multiplatform, Compose Multiplatform, Ktor, Navigation~3, OpenAPI Generator и др.);
		\item спроектировать архитектуру приложения на основе паттерна MVVM с использованием слоя репозиториев и ручного внедрения зависимостей;
		\item реализовать клиентское приложение с поддержкой аутентификации, управления пакетами достижений, отслеживания прогресса и интеграцией с серверным API;
		\item провести функциональное тестирование разработанного решения на целевых платформах.
	\end{enumerate}

	\textbf{Гипотеза исследования}: использование элементов геймификации (система достижений с поэтапным прогрессом, визуализация выполнения, обмен пакетами между пользователями) в кроссплатформенном приложении, реализованном на Kotlin Multiplatform, позволит создать эффективный инструмент мотивации и целеполагания, устраняющий выявленные недостатки существующих аналогов.

	\textbf{Практическая значимость} работы заключается в создании работоспособного кроссплатформенного приложения, которое может быть использовано для организации личных целей и отслеживания прогресса их достижения. Разработанное решение демонстрирует применение современных технологий Kotlin Multiplatform на практике и может служить основой для дальнейшего развития функциональности, включая офлайн-режим.

	\textbf{Публикации}: в рамках выполнения работы публикации не осуществлялись.

\end{document}
